\documentclass{beamer}

\usepackage{amsmath}
\usepackage{cancel}
\usepackage{booktabs}
\usetheme{metropolis}
\usepackage[style=authoryear]{biblatex}
\addbibresource{infillPuzzle.bib}
% Define Sauder green
\definecolor{saudergreen}{RGB}{91,191,33}
% Or: \definecolor{saudergreen}{HTML}{A3BD00}

% Apply it to the frame title background and text
\setbeamercolor{frametitle}{bg=saudergreen, fg=white}


\begin{document}

\title{Where does the infill housing go? \\Addition versus division}
\author{Thomas Davidoff, Robert Helsley, and Tsur Somerville\\Sauder, UBC}
\frame{\titlepage}

\begin{frame}
	\frametitle{Puzzle introduced and resolved}
	\begin{block}{Puzzle}
		\begin{itemize}
			\item Density limits have greatest deadweight loss in high-priced areas
				\begin{itemize}
					\item Across metros \textcite{HsiehMoretti}
					\item Within metros \textcite{bruecknerSingh}
				\end{itemize}
			\item But recent ``missing middle'' upzoning uptake - sorted within metros
		\end{itemize}
	\end{block}
	\begin{block}{Resolution}
		\begin{itemize}
			\item Addition of space, sure
			\item But a lot of infill upzoning is ``division''
				\begin{itemize}
					\item Same space, more units
					\item More valuable where poorer people live by concavity
					\item With mandated home sizes, rich get better locations
				\end{itemize}
		\end{itemize}
	\end{block}
\end{frame}

\begin{frame}
	\frametitle{Some examples from Vancouver: Dunbar  vs less rich Hastings-Sunrise}
	\pgfimage[width=1.2\textwidth]{dunbarHastings}
\end{frame}

\begin{frame}
	\frametitle{Literature}
	\begin{itemize}
		\item Ma: high-income households less harmed by large home mandates
		\item Studies showing negative sorting of infill takeup:
			\begin{itemize}
				\item Vancouver laneway (Davidoff, Pavlov, Somerville, 2022)
				\item ADUs in California (Marantz-Elmendorf and Kim, 2023; Brueckner and Thomaz, 2024)
				\item Duplexes and triplexes in Seattle (Huang et al, 2023)
				\item Portland's Residential Infill Project (RIP) (Dong and Hansz, 2019)
			\end{itemize}
	\end{itemize}
\end{frame}

\begin{frame}
	\frametitle{Spatially smoothed price per square foot Vancouver}
	 Price per square foot of single family homes in 2018, just before the reform, across locations in Vancouver. Top panel, spatially weighted transaction residuals 2014-2018, bottom panel lot-level BC Assessment estimates of land value per lot square foot
	\pgfimage[width=.4\textwidth]{surface_LocalPPSF.png}
	\pgfimage[width=.4\textwidth]{surface_bcaLogLandPPSF.png}
\end{frame}

\begin{frame}
	\frametitle{Elasticity spatial vancouver}
\end{frame}

\end{document}

\begin{frame}
	\frametitle{Overview}
	\begin{itemize}
		\item Introduce a puzzle regarding upzoning for ``infill housing''
			\begin{itemize}
				\item New density allowances in ``single family'' zones  \emph{should} be most valuable in high-priced areas
				\item But takeup of infill options leans towards lower-priced areas
			\end{itemize}
		\item Resolve the puzzle with a model
			\begin{itemize}
				\item ``addition'' (more square feet) vs ``division'' (more units)
				\item Land quality $\Rightarrow$ wealth $\Rightarrow$ less gain to division
				\item A lot of infill upzoning is division vs addition
			\end{itemize}
		\item Data: upzonings in Vancouver, Minneapolis, \cancel{Portland}
		\item Empirical contributions:
			\begin{itemize}
				\item Model assumption is true of single family zones:
					\begin{itemize}
						\item Price/sqft highly correlated with elasticity of price wrt sqft
					\end{itemize}
				\item Further evidence of negative sorting in division-only upzonings
				\item Some evidence of more positive sorting in mixed addition/division upzonings
			\end{itemize}
	\end{itemize}
\end{frame}

\begin{frame}
	\frametitle{Single vs multifamily zoning a very old question\\Brooks (1981) imagining Roman Senate}
	\begin{quote}
		Leader of Senate: All fellow members of the Roman senate hear me. Shall we continue to build palace after palace for the rich? Or shall we aspire to a more noble purpose and build decent housing for the poor? How does the senate vote?\\
		Entire Senate: F--- THE POOR!
	\end{quote}
\end{frame}

\begin{frame}
	\frametitle{Zoning deadweight loss greatest in best locations}
	\begin{itemize}
		\item Glaeser and Gyourko (2003) and Brueckner and Singh (2008): compare marginal price to cost /sqft
		\item Cities' natural shape: tallest, densest use closest to downtown
		\item For structure MC = price, need a very tall building downtown
		\item For land and structure  bang-for-buck equality,\\ need low density in suburbs
		\item Policy guidance re: upzoning to improve affordability: \\\begin{tiny} The model in this report identifies the lots most ripe for redevelopment under
the specified density scenarios, where the gains for converting from the exist-
ing single family properties are highest, as the size of the wedge. Consistent
with the basic model in urban economics, these are areas where land values
for single family lots are highest. In Metro Vancouver, the City of Vancouver
(especially its west side), the North Shore, Richmond and Burnaby stand
out. These results correspond well with expectations that high amenity ar-
eas closest in to the metropolitan centre would receive the greatest attention
for higher density redevelopment if it were enabled.\end{tiny}
\end{itemize}
\end{frame}
 
\begin{frame}
	\frametitle{Thus when Vancouver allowed 2 homes vs 1 per lot\dots Oh no!}
	\pgfimage[width=.7\textwidth]{duplexMap.png}
\end{frame}

\begin{frame}
	\frametitle{Not just Vancouver duplex where infill goes to \emph{low} $\theta$}
	\begin{itemize}
		\item Vancouver laneway (Davidoff, Pavlov, Somerville, 2022)
		\item ADUs in California (Marantz-Elmendorf and Kim, 2023; Brueckner and Thomaz, 2024)
		\item Duplexes and triplexes in Seattle (Huang et al, 2023)
		\item Portland's Residential Infill Project (RIP) (Dong and Hansz, 2019)
		\item Maybe we need a better model of infill location
	\end{itemize}
\end{frame}

\begin{frame}
	\frametitle{A model of infill location}
	\begin{itemize}
		\item ``primitives''
			\begin{itemize}
				\item $p(s,\theta)$: price per square foot of structure
					\begin{itemize}
						\item Size $s = \frac{f}{n}$
						\item $f$: floor area
						\item $n$: number of units
						\item Location $\theta$
						\item Lot size fixed at 1 unit of land
						\item Assume: $p_{\theta}>0$.
						\item Reduced form sorting (single family!): $p_{s\theta}>0$ 
					\end{itemize}
				\item Profit from raw land development $\pi = f\left[p(s,\theta) - c\right]$
					\begin{itemize}
						\item $c$: cost per sqft of structure, constant for now
						\item $f$ not necessarily freely chosen
					\end{itemize}
			\end{itemize}
	\end{itemize}
\end{frame}

\begin{frame}
	\frametitle{Positive sorting in theory?}
	\begin{itemize}
		\item Standard urban model, for each type, $$p_{\theta} \approx \frac{\text{Willingness to pay for }\theta}{\text{sqft}}$$
		\item Tradeoff: rich guys with more numerator also want denominator
		\item Can have poor people downtown apartments and rich people suburban detached homes
		\item But strict zoning deletes the denominator effect
		\item So only numerator $\Rightarrow$ expect + sorting
	\end{itemize}
\end{frame}


\begin{frame}
	\frametitle{Upzoning comparative statics addition versus division}
	\begin{itemize}
		\item Addition: allow extra $f$
			\begin{itemize}
				\item $\frac{\partial \pi}{\partial f} = p - c + sp_{s}$ 
				\item $\frac{\partial^{2}\pi}{\partial f\partial \theta} = p_{\theta}+sp_{s\theta} > 0$: more profitable in better locations
			\end{itemize}
		\item Division: allow more $n$
			\begin{itemize}
				\item $\frac{\partial \pi}{\partial n} = -s^{2}p_{s}$
				\item $\frac{\partial^{2}\pi}{\partial n \partial \theta} = -s^{2}p_{s\theta} < 0$: more profitable in worse locations
			\end{itemize}
	\end{itemize}
\end{frame}

\begin{frame}
	\frametitle{Upzonings in Vancouver}
	\begin{itemize}
		\item Predominant Zoning: R1-1 now, used to be several R zones
		\item Baseline was single family detached homes
		\item 1980s-2004: Basement suites permitted citywide 
		\item 2009: laneway homes permitted in single family homes
			\begin{itemize}
				\item Add square feet, but must be rented, lose garage space
				\item Laneway allowed .16-.25 FSR
			\end{itemize}
		\item Sept 2018: duplexes allowable throughout R zones
			\begin{itemize}
				\item No extra square feet (FSR .6-.7), but can sell 2 vs 1 home
				\item Generally \emph{can't} add a laneway
				\item Hence division and \emph{negative} addition 
			\end{itemize}
		\item 2023: multiplexes allowed
			\begin{itemize}
				\item 4-6 units permitted depending on lot size
					\begin{itemize}
						\item Much more attractive on wider lots
						\item Roughly 50\% of Vancouver lots 33' standard
					\end{itemize}
				\item FSR goes from .7 to 1.0
				\item Hence more division, more addition than duplex
				\item Also: single family down to .6 FSR, .7 for duplex
			\end{itemize}
	\end{itemize}
\end{frame}

\begin{frame}
	\frametitle{Data Sources}
	\begin{footnotesize}
	\begin{tabular}{llll}
	\hline
	City & Subject & Data Source & Coverage/Use \\
	\hline \hline
	Vancouver & Building Permits by  & City of Vancouver & 2019-2025 (laneway and duplex era)\\
	& &  & 2024-2025 (laneway, duplex, and multiplex era)\\
	Vancouver & Zoning & City of Vancouver & Current R1-1 zones  \\
	Vancouver & Transactions & B.C. Assessment & 2017-2019 for baseline price and elasticities \\
	Vancouver & Home and lot Characteristics & B.C. Assessment & All \\
	\hline
	Minneapolis & Building permit  & City of Minneapolis & 2020-2025 (post-policy)\\
	Minneapolis & Zoning & City of Minneapolis & Interior 1,2, and 3 zones\\
	Minneapolis & Transactions & Attom & 2016-2020 \\
	Minneapolis & Home and lot Characteristics & Attom\\
	\hline
	\end{tabular}
	\end{footnotesize}
\end{frame}


\begin{frame}
\frametitle{Is it true that $p_{s\theta}>0$ ?}
	\begin{itemize}
		\item BCA Single Family sales 2023-2024, Metro Vancouver
		\item Observation: a neighbourhood
		\item Elasticity $<-1$ means $p_{s}<0$, $>0$ means want lot assembly
			\begin{itemize}
				\item Both implausible
			\end{itemize}
	\end{itemize}
	\pgfimage[width=.8\textwidth]{ppsfElasticitysingleMetroRecent}
\end{frame}

\begin{frame}
	\frametitle{Price per square foot and elasticity: duplexs}
	\pgfimage[width=.8\textwidth]{ppsfElasticityduplexMetroRecent}
\end{frame}

\begin{frame}
\frametitle{Price per square foot and elasticity: condos\\Constant elasticity in size not a good fit\\Supply matters, too (competition should make close to 0)}
	\pgfimage[width=.8\textwidth]{ppsfElasticitycondoMetroRecent}
\end{frame}

\begin{frame}
	\frametitle{Spec vs custom: single vs not}
	\begin{itemize}
		\item Spec(ulative): developer buys, builds, sells to market
		\item Custom: rich guy buys, hires builder, holds property
		\item Hard to do custom other than single family
		\item Exercise: match CoV new build permit to subsequent sale, 2017-2021:
			\begin{itemize}
				\item Duplex 41 no, 77 yes
				\item Single 1884 no,706 yes
			\end{itemize}
		\item Custom likely fancier and not observed in sales data
		\item Elasticity hence likely biased down, more so West Side
		\item Note: can't do recent time slope of custom share: mechanical
	\end{itemize}
\end{frame}

\begin{frame}
	\frametitle{Not surprisingly, $p^{condo}_{s\theta}<p^{duplex}_{s/theta}<p^{sf}_{s\theta}$}
	\begin{itemize}
		\item Condos built to maximize $p-c$ 
		\item Singles too big
		\item Duplexes presumably ease waste of space
		\item Elasticity $<-1$ means $p_{s}<0$, $>0$ means want lot assembly
			\begin{itemize}
				\item Both implausible
			\end{itemize}
	\end{itemize}
\end{frame}


\begin{frame}
	\frametitle{City of Vancouver unit regressions by type}
	
\begingroup
\centering
\begin{tabular}{lccc}
   \tabularnewline \midrule \midrule
   Dependent Variable: & \multicolumn{3}{c}{logppsf}\\
   Model:         & (1)            & (2)            & (3)\\  
   \midrule
   \emph{Variables}\\
   logsqft        & -0.571$^{***}$ & -0.651$^{***}$ & -0.079$^{***}$\\   
                  & (0.002)        & (0.012)        & (0.002)\\   
   \midrule
   \emph{Fixed-effects}\\
   year           & Yes            & Yes            & Yes\\  
   NEIGHBOURHOOD  & Yes            & Yes            & Yes\\  
   \midrule
   \emph{Fit statistics}\\
   Observations   & 99,761         & 9,826          & 201,358\\  
   R$^2$          & 0.77528        & 0.58459        & 0.69394\\  
   Adjusted R$^2$ & 0.77480        & 0.57858        & 0.69374\\  
   \midrule \midrule
   \multicolumn{4}{l}{\emph{IID standard-errors in parentheses}}\\
   \multicolumn{4}{l}{\emph{Signif. Codes: ***: 0.01, **: 0.05, *: 0.1}}\\
\end{tabular}
\par\endgroup



\end{frame}

\begin{frame}
	\frametitle{Note dramatic price compression over time (foreign/spec?)\\Dot is a neighbourhood x year}
	\pgfimage[width=.7\textwidth]{logppsf_distribution_Van_single.png}
\end{frame}

\begin{frame}
	\frametitle{East/West Side and infill takeup over time: $\approx$ 33' lots}
	\pgfimage[width=.85\textwidth]{permitSharesByYearAndLotWidth33.png}
\end{frame}

\begin{frame}
	\frametitle{East/West Side and infill takeup over time: 50'+ lots}
	\pgfimage[width=.85\textwidth]{permitSharesByYearAndLotWidth50.png}
\end{frame}


\begin{frame}
	\frametitle{Neighbouhood x year Correlations: 20+, 33+ 40+ obs}
	\begin{footnotesize}
	\begin{tabular}{rrrrr}
  \hline
 & elasticity & intercept & duplex\_premium & shareDuplex \\
  \hline
elasticity & 1.00 & -0.99 & -0.89 & -0.11 \\
  intercept & -0.99 & 1.00 & 0.94 & 0.08 \\
  duplex\_premium & -0.89 & 0.94 & 1.00 & 0.01 \\
  shareDuplex & -0.11 & 0.08 & 0.01 & 1.00 \\
   \hline
   \hline
 & elasticity & intercept & duplex\_premium & shareDuplex \\
  \hline
elasticity & 1.00 & -0.99 & -0.84 & -0.15 \\
  intercept & -0.99 & 1.00 & 0.91 & 0.16 \\
  duplex\_premium & -0.84 & 0.91 & 1.00 & 0.14 \\
  shareDuplex & -0.15 & 0.16 & 0.14 & 1.00 \\
   \hline
  \hline
 & elasticity & intercept & duplex\_premium & shareDuplex \\
  \hline
elasticity & 1.00 & -0.99 & -0.91 & -0.22 \\
  intercept & -0.99 & 1.00 & 0.95 & 0.21 \\
  duplex\_premium & -0.91 & 0.95 & 1.00 & 0.17 \\
  shareDuplex & -0.22 & 0.21 & 0.17 & 1.00 \\
   \hline
\end{tabular}
\end{footnotesize}
\end{frame}

\begin{frame}
	\frametitle{Duplex vs single prices by East/Lot Width/Pre Post\\Percentage vs dollar difference, 2Duplex vs 1 single}
	\begin{footnotesize}
	\begin{tabular}{llllr}
	  \hline
	  east & widthCat & prePost & use & V1 \\ 
	  \hline
	    FALSE & wide & pre & single & 3746968.54 \\ 
	    FALSE & wide & pre & duplex & 2330138.76 \\ 
	   TRUE & wide & pre & single & 1303430.99 \\ 
	    TRUE & wide & pre & duplex & 1347562.60 \\ 
	  FALSE & wide & post & single & 4766580.72 \\ 
	    FALSE & wide & post & duplex & 2786995.62 \\ 
	    TRUE & wide & post & single & 2018256.53 \\ 
	    TRUE & wide & post & duplex & 1884056.27 \\ 
	    FALSE & standard & pre & single & 2684545.86 \\ 
	    FALSE & standard & pre & duplex & 1752426.75 \\ 
	    TRUE & standard & pre & single & 1558812.59 \\ 
	    TRUE & standard & pre & duplex & 1264738.68 \\ 
	    FALSE & standard & post & single & 3357192.98 \\ 
	    FALSE & standard & post & duplex & 2159571.82 \\ 
	    TRUE & standard & post & single & 2125509.22 \\ 
	    TRUE & standard & post & duplex & 1562854.13 \\ 
	   \hline
	\end{tabular}
  \end{footnotesize}
\end{frame}

\begin{frame}
	\frametitle{Minneapolis, MN}
	\begin{itemize}
		\item 2018 Council vote, effective 2020
		\item 3 regions: Interior-1, Interior-2, Interior-3
		\item All newly permit duplex and triplex vs SF only
		\item Interior-1: FSR (unch) 0.5 , max height 2.5 storeys (28 ft)
		\item Interior-2: FSR (unch) 0.5 , max height 2.5 storeys (28 ft)
		\item Interior-3: FSR up: 0.6 for duplexes, 0.7 for triplexes, max height 3 storeys (42 ft)
	\end{itemize}
\end{frame}

\begin{frame}
	\frametitle{Minneapolis Attom and permit data analysis\\Post-reform period only}
	\begin{itemize}
		\item For each zip code compute a mean price
		\item And an elasticity\\
		$\ln \frac{P_{izt}}{\text{sqft}_{izt}} = \beta_{0z} + \beta_{1z} \ln \text{sqft}_{izt} + \gamma X_{izt} + \delta_t + \epsilon_{izt}$
		\item These statistics computed 2 years pre-refom
		\item With permit data, compute fraction by zip code of permits
			\begin{itemize}
				\item Duplex or triplex vs single family
				\item In post-reform period only
			\end{itemize}
		\item Study correlations
		\item Expect: Interior-1 and -2: division only - sorting
		\item Interior-3: addition and division - more positive sorting
		\item Only 31 zips total (11 in UN-3)
	\end{itemize}
\end{frame}

\begin{frame}
	\frametitle{Sorting of infill in Minneapolis}
	\begin{footnotesize}
	\textbf{All affected Zips}
	\begin{tabular}{rrrrr}
	  \hline
	 & elasticity & price\_level & propensity & medianIncome \\ 
	  \hline
	elasticity & 1.00 & 0.59 & 0.28 & 0.30 \\ 
	  price\_level & 0.59 & 1.00 & 0.03 & 0.69 \\ 
	  propensity & 0.28 & 0.03 & 1.00 & -0.14 \\ 
	  medianIncome & 0.30 & 0.69 & -0.14 & 1.00 \\ 
	   \hline
	\end{tabular}
\textbf{Interior-3 zone}
\begin{tabular}{rrrrr}
  \hline
 & elasticity & price\_level & propensity & medianIncome \\ 
  \hline
elasticity & 1.00 & 0.61 & 0.32 & 0.31 \\ 
  price\_level & 0.61 & 1.00 & 0.24 & 0.82 \\ 
  propensity & 0.32 & 0.24 & 1.00 & 0.01 \\ 
  medianIncome & 0.31 & 0.82 & 0.01 & 1.00 \\ 
   \hline
\end{tabular}
\textbf{Interior-1 and 2 zones}
\begin{tabular}{rrrrr}
  \hline
 & elasticity & price\_level & propensity & medianIncome \\ 
  \hline
elasticity & 1.00 & 0.58 & 0.36 & 0.29 \\ 
  price\_level & 0.58 & 1.00 & -0.05 & 0.64 \\ 
  propensity & 0.36 & -0.05 & 1.00 & -0.30 \\ 
  medianIncome & 0.29 & 0.64 & -0.30 & 1.00 \\ 
   \hline
\end{tabular}
\end{footnotesize}
\end{frame}

\begin{frame}
	\frametitle{Conclusions}
	\begin{itemize}
		\item Infill zoning often more ``division'' than ``addition''
		\item Fact + model + assumption $\Rightarrow$ negative sorting
		\item The assumption is true: 
			\begin{itemize}
				\item More expensive \$psqft more $\frac{\partial \log p}{\partial \log s}$
			\end{itemize}
		\item Empirics: repeat usual negative sorting in:
			\begin{itemize}
				\item Vancouver duplex in division-only era
				\item Minneapolis division-only zones
			\end{itemize}
		\item Some evidence with small samples of more + sorting with addition
			\begin{itemize}
				\item Latter-day duplex if no laneway
				\item Multiplex 2025
				\item Interior-3 in Minneapolis
			\end{itemize}
		\item Our wish: figure a way to get more power in analysis of small N areas
	\end{itemize}
\end{frame}

\end{document}
